% ---------------------------------------------------------------------------
% R20/Reb2/S21: faulty RD voltage read-back (~54.2 V, top of the +/-55 V ADC
% range), which blocked the power-on of all three CCDs on the REB, and the
% software exception (LSSTCCSRAFTS-777) that recovered them.
% Sources: CCS log of the failed power-on of 2025-04-08 (Y. Utsumi);
% #cam-issues-ccd Slack discussion of 2025-04-08 to -17 (A. Roodman,
% S. Marshall, Y. Utsumi, S. Digel); Slack thread on the read-back hypothesis
% (S. Marshall, G. Thayer, C. Juramy) with REB schematic extracts and the REB
% board RD acceptance test; Slack transcript of the recovery session of
% 2025-04-24 (S. Marshall, S. MacBride, Y. Utsumi); SLAC JIRA
% LSSTCCSRAFTS-777; org-lsst-ccs-subsystem-rafts commit 829a8f2.
% No figures.
% ---------------------------------------------------------------------------

\subsubsection{R20/Reb2/S21: faulty RD voltage read-back}
\label{sec:r20reb2}

R20/Reb2 could not be powered on from the first days of April 2025 until
2025-04-24. Neither the CCDs nor the biases were at fault: the RD voltage
read-back of \texttt{Bias1}, the bias block of S21, had stuck at
$\simeq54.2$~V. That value lies in neither the ON nor the OFF window of the
bias level test, so the CCD power state of the board could not be resolved,
\texttt{setCCDsPowerState()} failed and the sequence was aborted
(\S\ref{sec:r20reb2-failure}). All three CCDs read by that REB, S20, S21 and
S22, were therefore off; they were the only CCDs in the focal plane not
powered on, and not exercised by the CCD shorts test, during the first week of
April.

The read-back is not a physically attainable voltage
(\S\ref{sec:r20reb2-notreal}), and the fault is understood as an open
connection between the RD sample point and the monitoring ADC, upstream of the
analogue multiplexer (\S\ref{sec:r20reb2-where}). The mitigation is a software
exception for this one channel \citep{jira:rafts777}: its RD read-back is
excluded from the shorts test, from the range check during power-on and from
the ON/OFF bias count, but only while the reading stays above 53~V
(\S\ref{sec:r20reb2-fix}). With that code in place the REB passed the shorts
test and powered on at the first attempt on 2025-04-24, and S21 delivers
nominal images (\S\ref{sec:r20reb2-recovery}). The cost is that S21 is
operated without an RD voltage monitor.

\paragraph{The failed power-on.}
\label{sec:r20reb2-failure}

\texttt{setCCDsPowerState()} evaluates eight tests on a REB: firmware bias
protection, the board-level ODI, CLKLI and CLKHI currents, the bias levels of
each of the three bias blocks, and the clock rails. On R20/Reb2 seven of them
passed. The attempt of 2025-04-08 reads

\begin{quote}\small\ttfamily
[2025-04-08T14:15:35.088+0000] SEVERE: R20/Reb2/Bias1 RD1 voltage = 54.21 NOT
in range for ON|OFF \\
$[$2025-04-08T14:15:35.089+0000] SEVERE: R20/Reb2/Bias1:
getBiasDacsPowerState() found 0(3) ON(OFF) of biases \\
$[$2025-04-08T14:15:35.089+0000] SEVERE: R20/Reb2: CCDs Bias voltage test
FAILED \\
$[$2025-04-08T14:15:35.089+0000] SEVERE: Test 6: FAILED Bias level test for
R20/Reb2/Bias1 \\
$[$2025-04-08T14:15:35.090+0000] SEVERE: R20/Reb2: setCCDsPowerState()
FAILED: 7:0:0 OFF:ON:DELTA of 8 tests passed \\
$[$2025-04-08T14:15:35.090+0000] SEVERE: R20/Reb2 CCDsPowerState updated to
FAULT
\end{quote}

\noindent followed by an automatic \texttt{powerCCDsOff()}. The OD, GD and OG
read-backs of that block were all in the OFF range, at $-0.03$, $-0.03$ and
$-0.01$~V, and \texttt{Bias0} and \texttt{Bias2} reported all four of their
biases OFF. \texttt{getBiasDacsPowerState()} required all four biases of a
block to agree before it would return OFF or ON, so with the RD read-back
outside both windows only three of four could ever be counted and the state
was unresolvable by construction, not merely out of tolerance. The board
stayed off throughout the period: \texttt{focal-plane/R20/Reb2/ODI} never
exceeded 8.8~mA.

\paragraph{The read-back is not a real voltage.}
\label{sec:r20reb2-notreal}

The channel returned only two distinct values, 54.18071 and 54.20757~V. Read
as adjacent codes of the 12-bit slow ADC, the step of
%
\begin{equation}
  54.20757 - 54.18071 = 0.02686~\mathrm{V}
  \label{eq:rdv-step}
\end{equation}
%
is one LSB, which implies a full-scale span of
$0.02686\times4096 = 110.0$~V, or $\pm55$~V. That is the range of this
channel. The REB driver scales the 12-bit conversions as

\begin{quote}\small\ttfamily
VOLT\_12\_SCALE\_R5 = 10.0 / 4095, \quad // +-5V ADC input range \\
VOLT\_12\_SCALE\_M5 = VOLT\_12\_SCALE\_R5 * 3, \\
VOLT\_12\_SCALE\_H5 = VOLT\_12\_SCALE\_R5 * 11,
\end{quote}

\noindent and \texttt{CHAN\_RD\_2} uses \texttt{VOLT\_12\_SCALE\_H5}. The ADC
input range is $\pm5$~V, so one code is $10/4095 = 2.442$~mV at the ADC input;
RD reaches that input through a divide-by-eleven network, so in units of RD
volts one code is $2.442\times11 = 26.862$~mV, matching
Eq.~(\ref{eq:rdv-step}) to five digits. The channel is therefore pinned near
the top of its range with about one count of noise, which is what a floating
ADC input looks like.

RD itself cannot reach that voltage. It is generated by a DAC output of at most
5~V multiplied by five, so it is limited to $\simeq25$~V, as verified in the
REB board acceptance tests; even a failure driving the output to its amplifier
supply would give the regulated 35~V, not 55~V. A reading of 54.2~V thus
cannot represent RD, and the bias itself was never suspect.

\paragraph{Where the fault is.}
\label{sec:r20reb2-where}

RD is monitored through the REB slow ADC, which multiplexes of order seventy
parameters. Following the board schematic, \texttt{VDD\_RD} is generated on
p.~32, divided on p.~31, enters an eight-channel analogue multiplexer on
p.~8, is buffered on p.~65 and digitised on p.~64. Other quantities on the
same multiplexer continued to read correctly, which places the fault upstream
of it, in the path specific to this channel. Since the failure is most likely
an open connection rather than a failed component, the behaviour may be
temperature dependent.

Three questions were left open. Whether the break is before the voltage
divider or after it was not established, the second case being the less
damaging. The raw ADC code was not recorded, and a value of identically zero,
as opposed to one moving by a few bits, would distinguish a divider tied to
ground from other failure modes; the corresponding channels of the other two
sensors on the board provide the comparison. No review of the affected
section of the board by an electronics engineer is documented.

\paragraph{Software mitigation.}
\label{sec:r20reb2-fix}

The decision was to make an exception for this channel in the CCD turn-on,
turn-off and status procedures and to ignore its value
\citep{jira:rafts777}. The exception is keyed on the bias path
\texttt{R20/Reb2/Bias1} together with the ADC identifier
\texttt{REBDevice.ADC\_RD\_1}, and is applied in three places in
\texttt{BiasControl.java} \citep{ccs:rafts777merge}:

\begin{itemize}
  \item \texttt{testShorts()} no longer drives and tests the RD channel of
        that block. The GD and OG channels are still tested. The skip is
        logged, so that it is visible in the CCS log.
  \item \texttt{checkAdc()}, which compares each read-back against its
        setpoint during power-on, skips the comparison for that channel
        \emph{only while the reading exceeds 53~V}. The exception is therefore
        self-disarming: should the read-back ever recover, the normal check
        applies again without a code change.
  \item \texttt{getBiasDacsPowerState()} counts three biases instead of four
        for that block, under the same 53~V condition, through a new variable
        \texttt{nbiases} that replaces the literal 4 in the OFF and ON
        conditions. This is the change that removes the \texttt{FAULT}.
\end{itemize}

The work was tracked as \texttt{LSSTCCSRAFTS-777}, raised on 2025-04-19,
reopened on 2025-05-05 as not yet merged, and resolved on 2025-05-07 with the
merge of commit \texttt{829a8f2}, released in
\texttt{org-lsst-ccs-subsystem-rafts} 3.0.10. The code was exercised on the
summit two weeks before that merge, as described next. The same commit also
carries changes to the OD ramp and to the sampling of the slow currents during
power-on that are unrelated to this mitigation.

\paragraph{Verification and recovery.}
\label{sec:r20reb2-recovery}

On 2025-04-24 the REB initialised cleanly with the patched CCS: firmware
\texttt{0x3139500e}, serial number 418254201, sequencer idle, six ASPICs
loaded, back bias enabled, \texttt{CCDsPowerState} OFF. The added log message
is not printed on a successful initialisation, so at that point the log did
not by itself confirm that the exception was in the running code.

A CCD shorts test was run on this REB alone before any attempt to power on:

\begin{quote}\small\ttfamily
[2025-04-24T19:57:38.742+0000] INFO: R20/Reb2 lowering sequencer clocks for
shorts test: 0x3bc --> 0xc \\
$[$2025-04-24T19:57:40.249+0000] INFO: R20/Reb2/Bias1: skipping testShorts()
on RD voltage, jira:LSSTCCSRAFTS-777 \\
$[$2025-04-24T19:57:44.258+0000] INFO: R20/Reb2 restoring sequencer clocks
after shorts test: 0xc --> 0x3bc \\
$[$2025-04-24T19:57:44.556+0000] INFO: testCCDShorts() finished successfully
\end{quote}

\noindent which confirmed both that the exception was active and that it
behaves as intended. In the trending of the test the block shows eight bias
excursions instead of the usual nine, three GD, two RD and three OG; the
missing one is S21/RD, whose read-back stays at $\simeq55$~V and dominates the
vertical scale unless that trace is excluded.

\texttt{powerCCDsOn} was then invoked once, with the agreement that work would
stop on any failure, and succeeded at the first attempt.
\texttt{focal-plane/R20/Reb2/ODI} read 91.3~mA and the per-CCD currents and
voltages were nominal. The faulty read-back moved slightly upward as the
biases came on, so the broken channel is not entirely decoupled from the ADC.
Clears and images followed, with stray light visible in the frames. A quick
statistics pass on the first image of S21,
\texttt{MC\_C\_20250424\_000016\_R20\_S21}, gives read noise of
3.44--3.70~ADU across the sixteen segments, serial CTE of 0.99997 or better
and parallel CTE of unity to the precision quoted, with no dead or anomalous
segment.

\paragraph{Operational consequences.}

Three CCDs were recovered by a change of a few tens of lines, and no hardware
intervention was needed or attempted. Two points follow from how it is done.
First, S21 runs without an RD voltage monitor: a genuine RD fault on that
sensor would no longer be caught by the bias level test, and the ON/OFF state
of that block is decided by three biases rather than four. Second, the
exception is hard-coded against a bias path string and a threshold rather than
being configuration, so it must be carried forward through releases of
\texttt{org-lsst-ccs-subsystem-rafts} and would silently stop protecting the
board if the path naming changed. Both should be tracked in configuration
management. The hardware fault is unrepaired and only partly characterised,
and the initialisation logging that would positively confirm that the
exception is active was left to be revised.
